% QR-CODE: Minitest Gleichstrommotor (Kl. 9, DS 09a)

\documentclass[a4paper,11pt]{article}

\usepackage[utf8]{inputenc}
\usepackage[T1]{fontenc}
\usepackage{helvet}
\renewcommand{\familydefault}{\sfdefault}

\usepackage[margin=20mm]{geometry}
\usepackage{tikz}
\usepackage{xcolor}
\usepackage{qrcode}
\usepackage{hyperref}
\hypersetup{colorlinks=true, urlcolor=physik}

\definecolor{physik}{HTML}{2c5aa0}
\colorlet{fachfarbe}{physik}

\newcommand{\lptitel}{Der Gleichstrommotor}
\newcommand{\lpurl}{https://devbox42.github.io/sciencesim/physik/kl09/motor-induktion/MT-09a-gleichstrommotor.html}
\newcommand{\lpurlkurz}{devbox42.github.io/.../MT-09a-gleichstrommotor.html}

\pagestyle{empty}

\begin{document}

\begin{center}

\vspace*{10mm}
{\fontsize{28}{34}\selectfont\bfseries\color{fachfarbe}Minitest: \lptitel}

\vspace{15mm}

\href{\lpurl}{\fcolorbox{fachfarbe}{white}{\qrcode[height=100mm]{\lpurl}}}

\vspace{12mm}

{\large\color{gray}\lpurlkurz}

\vspace{8mm}

{\Large\color{gray}Scanne den QR-Code mit deinem Gerät}

\vspace{10mm}

\fcolorbox{fachfarbe}{fachfarbe!10}{%
    \parbox{0.7\textwidth}{%
        \centering
        \vspace{3mm}
        {\large Stundenleistung: Gleichstrommotor}\\[2mm]
        6 Aufgaben -- 25 BE -- ca. 15 Minuten
        \vspace{3mm}
    }%
}

\end{center}

\end{document}
